\documentclass{standalone}

\newif\ifwhitecard
\whitecardtrue

\usepackage{microtype}
\usepackage{tikz}
\usepackage{xcolor}
\usepackage{xspace}

% fonts
\usepackage{fontspec}
% Using Linux-compatible fonts
\newfontfamily\comicsans{DejaVu Sans} % Alternative to Comic Sans
\setsansfont{Liberation Sans} % Alternative to Helvetica

% font settings for default text
\newcommand{\bodyfont}{\sffamily\bfseries\fontsize{12}{12}\selectfont}
\newcommand{\bookletfont}{\sffamily\fontsize{12}{12}\selectfont}
\newcommand{\titlefont}{\sffamily\bfseries\fontsize{20.5}{20.5}\selectfont}

% prevent hyphenation
\pretolerance=10000
\tolerance=2000 
\emergencystretch=10pt

% Define REFERENCE and OPTIONAL commands to handle citations
\newcommand{\REFERENCE}[1]{} % Empty command that takes one argument
\newcommand{\OPTIONAL}[1]{} % Empty command that takes one argument

% BLANK
\newcommand{\BLANK}{{\color{cardfg}\tikz[baseline] \draw [line width=1.5pt] (0,0)--+(0.5in,0);\nobreak\color{cardfg}}}

\ifwhitecard
	\definecolor{cardbg}{named}{white}
	\definecolor{cardfg}{named}{black}
	\definecolor{icondark}{named}{black}
	\definecolor{iconlight}{named}{white}
\else
	\definecolor{cardbg}{named}{black}
	\definecolor{cardfg}{named}{white}
	\definecolor{icondark}{named}{black}
	\definecolor{iconlight}{named}{white}
\fi

\newcommand{\cardicon}{
	\begin{tikzpicture}[yscale=-1]
	\draw [color=cardfg,fill=icondark,line width=0.5pt,rotate around={-20:(0.19in,0.105in)}] (0.075in,0.105in) rectangle +(0.19in,0.19in);
	\draw [color=icondark,fill=iconlight,line width=0.5pt] (0.19in,0.105in) rectangle +(0.19in,0.19in);
	\draw [color=icondark,fill=iconlight,line width=0.5pt,rotate around={20:(0.19in,0.105in)}] (0.275in,0.073in) rectangle +(0.19in,0.19in);
	\end{tikzpicture}
}

\input{booklet}

\begin{document}
\begin{tikzpicture}[yscale=-1]

\renewcommand{\baselinestretch}{0.8}

\drawbackground

% title page
\node [cardfg,below right] at (2.2in + 0.12in, 0.12in) {\titlefont Cards};
\node [cardfg,below right] at (2.2in + 0.12in, 0.42in) {\titlefont Against};
\node [cardfg,below right] at (2.2in + 0.12in, 0.72in) {\titlefont Blockchain};
\node [above right] at (2.2in + 0.12in, 3.4in - 0.12in) {\cardicon};

% back page
\node [cardfg,below right,text width=1.96in] at (0.12in, 0.12in) {
	{\sffamily\fontsize{8}{8}\selectfont
	Cards Against Blockchain is shamelessly based on Cards Against Cryptography, which was released under a Creative Commons BY-NC-SA 2.0 license.}

	\medskip

	{\sffamily\fontsize{8}{8}\selectfont
	Cards Against Blockchain is released under the same license, which means you can use, remix, and share the game for free, but cannot sell it.}

	\bigskip

	\includegraphics[width=0.25in]{icons/cc}
	\includegraphics[width=0.25in]{icons/by}
	\includegraphics[width=0.25in]{icons/nc}
	\includegraphics[width=0.25in]{icons/sa}

	\bigskip

	{\sffamily\fontsize{8}{8}\selectfont
	The source code and printable PDF versions of Cards Against Blockchain are available online at https://github.com/CardsAgainst\\Blockchain/. You can also submit pull requests to propose new cards.
	}

	\medskip

%	{\sffamily\fontsize{6}{6}\selectfont
%	Eurocrypt 2019 edition.
%	}

};
\end{tikzpicture}
\end{document}
